\documentclass[11pt,a4paper]{article}
\usepackage[utf8]{inputenc}
\usepackage[margin=0.85in]{geometry}
\usepackage{amsmath,amssymb,amsfonts}
\usepackage{graphicx}
\usepackage{booktabs}
\usepackage{xcolor}
\usepackage{fancyhdr}
\usepackage{titlesec}
\usepackage{float}
\usepackage{listings}
\usepackage{hyperref}

\definecolor{navyblue}{RGB}{0, 51, 102}
\definecolor{darkslate}{RGB}{47, 79, 79}
\definecolor{codebg}{RGB}{245, 245, 245}

\hypersetup{
    colorlinks=true,
    linkcolor=navyblue,
    citecolor=navyblue,
    urlcolor=navyblue,
    pdftitle={Empirical Task Methods: Probabilistic Reversal Continuous Bridge},
    pdfauthor={Lead Computational Architect}
}

\lstset{
    backgroundcolor=\color{codebg},
    basicstyle=\ttfamily\footnotesize,
    breaklines=true,
    frame=single,
    framerule=0.5pt,
    rulecolor=\color{darkslate}
}

\pagestyle{fancy}
\fancyhf{}
\rhead{\textcolor{darkslate}{\small Empirical Task Methods \& Continuous Injection Bridge}}
\lhead{\textcolor{darkslate}{\small Technical Methodology Report}}
\rfoot{\thepage}

\titleformat{\section}{\large\bfseries\color{navyblue}}{\thesection}{1em}{}
\titleformat{\subsection}{\bfseries\color{darkslate}}{\thesubsection}{1em}{}

\title{\Huge \textbf{\color{navyblue} Empirical Task Methodology \& Continuous Injection Mapping}\\[0.4em]
\large Formal Protocol Definition for 2AFC Probabilistic Reversal Data Ingestion}
\author{\textbf{Lead Computational Architect \& Antigravity Research Group}}
\date{August 15, 2026}

\begin{document}

\maketitle

\begin{abstract}
This technical methodology report defines the exact data dictionary, temporal structure, and discrete-to-continuous injection mapping for the 2AFC Probabilistic Reversal behavioral dataset (\texttt{behavioral\_compilate.csv}). Comprising 16,029 trials across 128 human participants, the task pairs explicit reward magnitudes ($M_A, M_B \in \{2..8\}$) with hidden probabilistic reward contingencies ($85/15$ vs $15/85$). We formulate a mathematically rigorous continuous-time Mossy Fiber injection bridge ($\mathbf{u}_t$) with $\Delta t = 10\text{ ms}$ resolution, enabling seamless deployment of the \texttt{ExactRModel} continuous-time C++ cerebellar reservoir.
\end{abstract}

\vspace{0.8em}
\hrule
\vspace{0.8em}

\section{Behavioral Task Design \& Experimental Structure}

The empirical dataset captures human decision-making in a 2-Alternative Forced Choice (2AFC) Probabilistic Reversal Learning paradigm:
\begin{enumerate}
    \item \textbf{Stimulus Presentation}: Participants are presented with two options ($A$ and $B$) displaying explicit numerical reward magnitudes ($M_A = \text{Bd1}, M_B = \text{Bd2} \in \{2, 3, 4, 5, 6, 7, 8\}$).
    \item \textbf{Hidden Probabilistic Contingency}: One option offers a high reward probability ($85\%$), while the other offers a low reward probability ($15\%$). Contingency states periodically reverse without warning.
    \item \textbf{Stochastic Feedback}: Upon choice selection, participants receive either the full magnitude ($M_A$ or $M_B$) if rewarded ($F = 1$), or $0$ points if unrewarded ($F = 0$).
\end{enumerate}

\section{Data Dictionary \& Summary Statistics}

The raw dataset \texttt{behavioral\_compilate.csv} contains 16,029 rows across 128 subjects.

\begin{table}[H]
\centering
\caption{Empirical Dataset Data Dictionary (\texttt{behavioral\_compilate.csv}).}
\label{tab:data_dict}
\vspace{0.5em}
\footnotesize
\begin{tabular}{l l l}
\toprule
\textbf{Variable} & \textbf{Type} & \textbf{Description \& Range} \\
\midrule
\texttt{participant\_id} & String & Unique subject identifier (e.g., \texttt{ACMO\_06011994}). \\
\texttt{nt} & Integer & Sequential trial index within subject ($1 \dots 130$). \\
\texttt{buena} & Binary & Latent state ($1 \implies A$ is $85\%$ optimal; $0 \implies B$ is $85\%$ optimal). \\
\texttt{Bd1} ($M_A$) & Integer & Explicit reward magnitude for Option A ($2 \dots 8$). \\
\texttt{Bd2} ($M_B$) & Integer & Explicit reward magnitude for Option B ($2 \dots 8$). \\
\texttt{Resp} & Categorical & Choice executed ($1 \implies \text{Option A}$; $2 \implies \text{Option B}$). \\
\texttt{F} & Binary & Feedback outcome ($1 \implies \text{Reward Delivered}$; $0 \implies \text{Zero Points}$). \\
\texttt{prob} & Continuous & True underlying probability for Option A ($0.15$ or $0.85$). \\
\texttt{ttp} & Float (ms) & Timestamp of Stimulus Presentation onset. \\
\texttt{ttr} & Float (ms) & Timestamp of Participant Response execution. \\
\texttt{ttF} & Float (ms) & Timestamp of Feedback Presentation onset. \\
\bottomrule
\end{tabular}
\end{table}

\subsection{Summary Statistics}
\begin{itemize}
    \item \textbf{Total Participants}: $128$
    \item \textbf{Total Observations}: $16,029$ trials (Mean $125.2$ trials/subject)
    \item \textbf{Overall Accuracy}: $49.86\%$ choice of optimal contingency option
    \item \textbf{Overall Reward Rate}: $63.42\%$
\end{itemize}

\section{The Continuous Injection Bridge ($\mathbf{u}_t$ Mapping)}

Because the C++ \texttt{ExactRModel} cerebellar reservoir operates in continuous time ($\Delta t = 10\text{ ms}$), discrete event timestamps ($\text{ttp} \to \text{ttr} \to \text{ttF}$) must be transformed into a 5-dimensional Mossy Fiber input vector $\mathbf{u}(t) \in \mathbb{R}^5$:
\begin{equation}
\mathbf{u}(t) = \left[ M_A(t), M_B(t), \text{Choice}_A(t), \text{Choice}_B(t), \text{Outcome}(t) \right]^T
\end{equation}

\subsection{Temporal Phase Array Padding Scheme}
For each trial $k$, the discrete interval $[t_0, t_f]$ is discretized into time steps $t = m \cdot \Delta t$:

\begin{enumerate}
    \item \textbf{Stimulus Window} ($t \in [\text{ttp}_k, \text{ttr}_k]$):
    $$M_A(t) = \text{Bd1}_k, \quad M_B(t) = \text{Bd2}_k, \quad \text{Choice}_{A,B}(t) = 0, \quad \text{Outcome}(t) = 0$$
    \item \textbf{Response Window} ($t \in [\text{ttr}_k, \text{ttF}_k]$):
    $$\text{Choice}_A(t) = \mathbb{I}(\text{Resp}_k = 1), \quad \text{Choice}_B(t) = \mathbb{I}(\text{Resp}_k = 2), \quad \text{Outcome}(t) = 0$$
    \item \textbf{Feedback Window} ($t \in [\text{ttF}_k, \text{ttF}_k + 200\text{ms}]$):
    $$\text{Outcome}(t) = F_k \cdot M_{\text{chosen}, k}$$
    \item \textbf{Inter-Trial Interval (ITI)} ($t \in [\text{ttF}_k + 200\text{ms}, \text{ttp}_{k+1}]$):
    $$\mathbf{u}(t) = \mathbf{0}$$
\end{enumerate}

This continuous bridge preserves exact sensorimotor dynamics, ensuring the Granule cell layer receives magnitude stimuli prior to decision execution.

\end{document}
